% =====================================================================
%  CONJURA - SELF-CONTAINED NOTE
%  The capping construction for c/0010, with the whole chain it rests on.
%  Requires conjura-solution.cls and conjura-conjecture.cls here.
% =====================================================================
\documentclass{conjura-solution}

\runninghead{CAPPING THE FIXED SET}

\newcommand{\Fun}{\mathsf{Fun}}
\newcommand{\Adv}{\mathsf{Adv}}
\newcommand{\Real}{\mathsf{Real}}
\newcommand{\Dec}{\mathsf{Dec}}
\newcommand{\Prob}[1]{\Pr\!\left[#1\right]}
\newcommand{\Exp}{\mathbb{E}}
\newcommand{\bits}{\{0,1\}}
\newcommand{\vx}{\mathbf{x}}
\newcommand{\vz}{\mathbf{z}}
\newcommand{\vu}{\mathbf{u}}
\newcommand{\ind}{\mathbb{1}}
\newcommand{\Unif}{\mathrm{Unif}}
\newcommand{\supp}{\mathrm{supp}}
\newcommand{\SD}{\mathrm{SD}}
\newcommand{\kna}{\kappa^{\mathrm{na}}}
\newcommand{\Ass}[1]{\textbf{[#1]}}

\newtheorem*{lemHz}{Lemma H\textlf{0}}
\newtheorem*{lemHo}{Lemma H\textlf{1}}
\newtheorem*{thmH}{Theorem H}
\newtheorem*{thmHp}{Theorem H$'$}
\newtheorem*{assumption}{Assumption}

\begin{document}

\cjkicker{CONJURA \textperiodcentered{} SELF-CONTAINED NOTE}
\cjtitle{Capping the Fixed Set}
\cjsubtitle{The $P$-axis of the split-source decomposition, with the whole chain it rests on}
\cjresolves{\emph{Decomposition for Split Unpredictable Sources}
  (\texttt{c/0010}). \textbf{Partially, and for a relaxed statement.} The
  conjecture's hypothesis on $P$ is removed entirely, on the region
  $(\mathrm{H2})$, provided the mixture family may depend on the observer's query
  budget $q$. That is a relaxation of the conjecture, not the conjecture. The
  order-respecting statement, with an absolute constant and a $q$-free family, is
  \emph{not} settled here.}

\cjstatus{Self-contained modulo four published external results, which are isolated
  as Assumptions in \S\ref{sec:external} and are the only unproved inputs.
  Everything else is stated and proved below. Verification is \emph{not} uniform
  across the chain, and \S\ref{sec:verification} gives it per result: the final
  arm has had six independent blind referee passes, the arm it rests on has had
  none.}

\vspace{0.4em}

\section*{Reader's guide}

\noindent This note proves, on a region, a relaxation of a conjecture about
presampling for two sources that read a random oracle independently. It is
assembled from four campaign artifacts so that it can be read without them.

\smallskip
\noindent\textbf{The new content is \S\ref{sec:capping}, and it is short.} The
conjecture asks for \emph{some} family of $P$-mixtures; it never asks that the
fixed sets be large, and Definition~\ref{def:bf} bounds them by $|I| \le P$ rather
than fixing $|I|$. So when the requested $P$ exceeds what the proof can afford, do
not use it --- build the family at the balance point instead. A $P_0$-mixture with
$P_0 \le P$ \emph{is} a $P$-mixture. That removes the hypothesis $(\mathrm{H1})$
from Theorem~\ref{thm:Epp} at a cost of $8 \to 13$ in the constant $C$.

\smallskip
\noindent\textbf{Everything before \S\ref{sec:capping} is the chain.}
\S\S\ref{sec:structural}--\ref{sec:lower} build two extraction bounds and one
lower bound; \S\ref{sec:presampling} constructs the family;
\S\ref{sec:decomposition} converts extraction into decomposition;
\S\ref{sec:reduce} reduces the hypotheses. A reader who wants only the new result
can go to \S\ref{sec:capping} and take Theorem~\ref{thm:Epp} on faith.

\smallskip
\noindent\textbf{Two things this does not prove, both overclaimed in an earlier
draft before referees caught them.}

\begin{itemize}[leftmargin=1.2em,itemsep=2pt]
\item The family must be allowed to depend on $q$, which
  Remark~\ref{rem:order} forbids. So Theorem~H proves a relaxation, and the
  window recorded as open --- absolute $C$, $q$-free family --- stays open.
\item Keeping the family $q$-free costs a factor $\sqrt{q^{+}}$ in the constant
  (Theorem~H$'$). It does not follow that the $\sqrt{q^{+}}$ \emph{is} the price of
  $q$-independence: what is shown is that $q$-dependence \emph{suffices} to remove
  it within this route. No lower bound against $q$-free families exists anywhere.
\end{itemize}

\smallskip
\noindent\textbf{What remains.} The hypothesis $(\mathrm{H2})$, a condition on the
output size $M$, is carried unchanged throughout. It is the whole residue on the
$P$-axis; \S\ref{sec:h2} is about what it means and whether applications satisfy
it. They do for key extraction and do not for long-output extraction.

\section{Setting}\label{sec:setting}

Fix $N, M \ge 2$ and let $\Fun := \{f : [N] \times [N] \to [M]\}$, with
$H \gets_{\$} \Fun$. Logarithms are to base $2$ and $\ln$ is natural.

\begin{definition}[sources, splitness, unpredictability]\label{def:sources}
A \emph{source} is a computationally unbounded randomized algorithm $S$ with
oracle access to $H$ and unbounded query count --- so it may read all of $H$ ---
taking no input, using private coins, and outputting a pair
$(x,z) \in [N] \times \bits^{*}$. Two sources $S_1, S_2$ are \emph{split} if run
on independent private coins, so that conditioned on $H$ their output pairs are
independent. Write $(x_i,z_i) \gets S_i^{H}$ with $|z_i| \le \sigma_i$, and put
$\vx := (x_1,x_2)$, $\vz := (z_1,z_2)$.

A \emph{predictor} $\mathsf{P}$ is computationally unbounded with oracle access to
$H$ and unbounded query count, takes $\vz$, and outputs an element of $[N]$. The
pair is \emph{$\delta$-unpredictable} if
$\Prob{\mathsf{P}^{H}(\vz) = x_i} \le \delta$ for $i = 1,2$ and every $\mathsf{P}$.
Since a predictor may ignore its input and output a fixed element, necessarily
$\delta \ge 1/N$.
\end{definition}

\begin{definition}[bit-fixing sources, and consistency]\label{def:bf}
A distribution $Y$ over $\Fun$ is \emph{$P$-bit-fixing} if there are a set
$I \subseteq [N]\times[N]$ with $|I| \le P$ and a map $a : I \to [M]$ such that a
draw $g \gets Y$ satisfies $g(u) = a(u)$ for every $u \in I$ and is uniform on
$[M]$ and independent across the points $u \notin I$. We call $I$ its \emph{fixed
set} and $a$ its \emph{fixed values}, and say $Y$ is \emph{consistent with}
$f \in \Fun$ if $a(u) = f(u)$ for every $u \in I$. A \emph{$P$-mixture} is a
finite indexed family $(\lambda_j, Y_j)_{j \in \mathcal{J}}$ with $\lambda_j \ge 0$,
$\sum_j \lambda_j = 1$ and each $Y_j$ a $P$-bit-fixing source; it is
\emph{consistent with} $f$ if every component is. A \emph{draw} means: draw $J$
with $\Prob{J = j} = \lambda_j$, then $g \gets Y_J$.
\end{definition}

The inequality $|I| \le P$ is the hinge of \S\ref{sec:capping} and is not a
stylistic choice: Assumption~\ref{ass:cdgs2}'s source states the same convention,
``$P$-bit-fixing if it is fixed on \emph{at most} $P$ coordinates''.

\begin{definition}[observers, and the four experiments]\label{def:obs}
An \emph{observer} is a computationally unbounded algorithm $D$ with oracle access,
making at most $q$ queries, taking a pair in $[M] \times (\bits^{*})^{2}$ and
outputting a bit. It is \emph{challenge-oblivious}: its input is $\vz$ together
with a single value of $[M]$, never $\vx$. Given a family
$\mathcal{Y} = \{Y_{f,\zeta}\}$ of $P$-mixtures indexed by $f \in \Fun$ and
$\zeta \in (\bits^{*})^{2}$:
\begin{itemize}[leftmargin=1.4em,itemsep=1pt]
\item $\Real$: draw $H$; draw $(x_i,z_i) \gets S_i^{H}$ on independent coins; set
  $y := H(\vx)$; output $D^{H}(y,\vz)$.
\item $\Real_0$: as $\Real$ but $y \gets_{\$} [M]$.
\item $\Dec$: draw $H$ and $(\vx,\vz)$ as in $\Real$ and do not resample; draw
  $(J,H^{*}) \gets Y_{H,\vz}$; set $y^{*} := H^{*}(\vx)$; output
  $D^{H^{*}}(y^{*},\vz)$. Write $I_J$ for the fixed set drawn.
\item $\Dec_0$: as $\Dec$ but $y \gets_{\$} [M]$.
\end{itemize}
Put $\Adv_{\mathcal{Y},D} := |\Prob{\Real = 1} - \Prob{\Dec = 1}|$ and
\[
  \kappa(q) := \sup_{D} \bigl|\Prob{\Real = 1} - \Prob{\Real_0 = 1}\bigr| ,
\]
the supremum over $q$-query challenge-oblivious $D$. Note $\kappa$ involves
$\Real$ and $\Real_0$ only, so it does not depend on $\mathcal{Y}$: the oracle is
the real $H$ in both and only the challenge changes. Write $\kna(q)$ for the same
supremum restricted to observers of challenge resolution $1$
(Definition~\ref{def:res}).
\end{definition}

Throughout,
\[
  \sigma := \sigma_1 + \sigma_2, \qquad
  \sigma' := \sigma + 2\log N, \qquad
  q^{+} := q + 1, \qquad
  S := \sqrt{\sigma' q^{+} \delta},
\]
\[
  t_q := \sqrt{\sigma' q^{+}/\delta}, \qquad
  t_0 := \sqrt{\sigma'/\delta}, \qquad
  t_q = t_0\sqrt{q^{+}}, \qquad
  B := \sigma' + \log \gamma^{-1},
\]
and $\mathcal{Z} := \bits^{\le \sigma_1} \times \bits^{\le \sigma_2}$, so
$|\mathcal{Z}| < 2^{\sigma+2}$. Since $N \ge 2$ gives $2 \log N \ge 2$ and
$\sigma \ge 0$, we have $\sigma' \ge 2$ throughout; this is used repeatedly and
is why $N \ge 2$ is a standing hypothesis.

\begin{conjecture}[the target]\label{conj:main}
There are absolute constants $c, C$ such that for all $N, M$, every split
$\delta$-unpredictable pair with leakage bounded by $\sigma_1,\sigma_2$, every
$P \in \mathbb{N}$ and every $\gamma \in (0,1)$, there is a family
$\mathcal{Y} = \{Y_{f,\zeta}\}$, $f \in \Fun$, $\zeta \in (\bits^{*})^{2}$,
depending only on $S_1, S_2, P, \gamma$, in which every $Y_{f,\zeta}$ is a
$P$-mixture consistent with $f$, such that for every $q \in \mathbb{N} \cup \{0\}$
and every $q$-query challenge-oblivious $D$,
\[
  \Adv_{\mathcal{Y},D} \le
  \frac{c\bigl(\sigma' + \log \gamma^{-1}\bigr) q^{+}}{P}
  + C\sqrt{\sigma' q^{+} \delta} + \gamma .
\]
\end{conjecture}

\begin{remark}[what the indexing allows]\label{rem:index}
The family is indexed by the oracle as well as the leakage --- the one departure
from the classical single-source lemma. It may \emph{not} be chosen using $\vx$,
which is why the index set is $\Fun \times (\bits^{*})^{2}$ and nothing larger.
\end{remark}

\begin{remark}[quantifier order]\label{rem:order}
The order is part of the statement. The family is chosen after $P$ and $\gamma$ but
\emph{before} $q$, so no part of it may be tuned to the observer's query budget.
Relaxing this makes the statement easier. Theorem~H below takes exactly that
relaxation and no other.
\end{remark}

Two facts about $\vx$ are used in \S\ref{sec:decomposition}. Both are the
Contract's, and neither is reproved here.

\begin{lemma}[the fixed set misses the challenge]\label{lem:hit}
Let $\mathcal{Y}$ be any family as in Conjecture~\ref{conj:main}, let $J$ index the
component drawn in $\Dec$ or $\Dec_0$, and let $I_J$ be its fixed set. Then
$\Prob{\vx \in I_J} \le P\delta$.
\end{lemma}

\begin{lemma}[the observer misses the challenge]\label{lem:query}
In $\Dec_0$, let $Q$ be the set of points $D$ queries. Then
$\Prob{\vx \in Q} \le q\delta$, the case $q = 0$ holding because $Q$ is empty.
\end{lemma}

\section{The external inputs}\label{sec:external}

Four published results are used and not proved. They are isolated here so that
the rest of the note is self-contained relative to them. Each is quoted from a
source card rather than from the paper itself, which is a gap recorded in
\S\ref{sec:verification}.

\begin{assumption}[\Ass{CDGS, Claim 2}, decomposition]\label{ass:cdgs2}
Let $\mu$ be a distribution on $\Fun$ with $S := N^{2}\log M - H_\infty(\mu)$, let
$\delta_0 \in (0,1]$ and $\gamma \in (0,1)$, and put
$P := (S + \log \gamma^{-1})/(\delta_0 \log M)$. Then
$\mu = \sum_j \lambda_j X_j + \lambda_{\mathrm{fin}} Y_{\mathrm{fin}}$ with
$\lambda_{\mathrm{fin}} \le \gamma$, the supports pairwise disjoint, and each $X_j$
a $(P, 1-\delta_0)$-dense source fixing a set $I_j$ with $|I_j| \le P$ to values
every $f$ in its support takes.
\end{assumption}

\begin{assumption}[\Ass{CDGS, Claim 3}, dense-versus-bit-fixing]\label{ass:cdgs3}
Let $X$ be a $(P,1-\delta_0)$-dense source and $Y$ the $P$-bit-fixing source on the
same fixed set and values. Every $q$-query adaptive distinguisher between $X$ and
$Y$ has advantage at most $q\delta_0 \log M$.
\end{assumption}

\begin{assumption}[\Ass{CFHS, Lemma 3}, flat decomposition]\label{ass:flat}
Let $k \in \mathbb{R}_{>0}$ with $2^{k} \in \mathbb{N}$. Every $k$-source $X$ ---
one with $\Prob{X = x} \le 2^{-k}$ for all $x$ --- is a convex combination of flat
$k$-sources, each uniform on a set of size exactly $2^{k}$.
\end{assumption}

\begin{assumption}[\Ass{CFHS, Lemma 4.3} and Hoeffding]\label{ass:binom}
For $k,n \in \mathbb{N}$ with $0 < k \le n$, $\ \binom{n}{k} \le (en/k)^{k}$. And
if $Y_1,\dots,Y_n$ are independent, mean zero, each in an interval of length $1$,
then $\Prob{|\sum_i Y_i| \ge \lambda} \le 2\exp(-2\lambda^{2}/n)$.
\end{assumption}

\section{Structural facts}\label{sec:structural}

\begin{lemma}[derandomisation]\label{lem:0}
Fix one of the two paired experiments. For every $q$ and every $q$-query
challenge-oblivious $D$ there is a deterministic such $D'$ with the relevant
advantage no smaller for $D'$ than for $D$, and of the same challenge resolution.
Consequently the suprema defining $\kappa(q)$ and $\kna(q)$ are attained on
deterministic observers, the latter on deterministic observers of resolution $1$.
\end{lemma}

\begin{proof}
Let $\rho$ be $D$'s coins and $D_\rho$ the deterministic strategy it then runs. All
of $\Fun, [N], [M], \mathcal{Z}$ are finite and $D$ makes at most $q$ queries, so
there are finitely many deterministic $q$-query challenge-oblivious strategies. The
quantity at issue is $|\Exp_\rho[\alpha(\rho)]|$ with
$\alpha(\rho) = \Prob{\mathsf{E}_1 = 1 \mid \rho} - \Prob{\mathsf{E}_0 = 1 \mid \rho}$,
the coins being independent of every other draw and the paired experiments
differing in no other component; so it is at most $\max_\rho |\alpha(\rho)|$,
attained at some $\rho^{*}$. Take $D' := D_{\rho^{*}}$. It inherits $D$'s
resolution because Definition~\ref{def:res} quantifies over every coin string, so a
property holding at every $\rho$ holds at $\rho^{*}$.
\end{proof}

Once $D$ is deterministic and its oracle fixed to $f$, its output is a function of
its challenge input alone; write $\theta : [M] \to \bits$ for that function and
$p_\theta := |\theta^{-1}(1)|/M$.

\begin{lemma}[posterior product structure]\label{lem:1}
For every $(f,\zeta)$ in the support of $(H,\vz)$, the conditional law of $\vx$ is
the product $\Pi_{f,\zeta} = \pi^{1}_{f,\zeta_1} \otimes \pi^{2}_{f,\zeta_2}$ with
$\pi^{i}_{f,\zeta_i}(u) = \Prob{x_i = u \mid H = f, z_i = \zeta_i}$. Writing
$m_i := \|\pi^{i}_{f,\zeta_i}\|_\infty$:
(i) $\max_{\vu} \Pi_{f,\zeta}(\vu) = m_1m_2$;
(ii) $\Exp[m_i] \le \delta$;
(iii) $\Exp[m_1m_2] \le \delta$;
(iv) $\Prob{m_1m_2 > \vartheta} \le 2\delta/\sqrt{\vartheta}$ for
$\vartheta \in (0,1]$.
\end{lemma}

\begin{proof}
The sources run on independent private coins and on the same $H$, so conditioned on
$H = f$ their output pairs are independent:
$\Prob{x_1 = u_1, z_1 = \zeta_1, x_2 = u_2, z_2 = \zeta_2 \mid f}
 = \prod_i \Prob{x_i = u_i, z_i = \zeta_i \mid f}$.
Summing over $u_1,u_2$ gives
$\Prob{\vz = \zeta \mid f} = \prod_i \Prob{z_i = \zeta_i \mid f} > 0$, and dividing
gives the product form. (i) is that the maximum of a product measure is the product
of the maxima.

For (ii), let $\mathsf{P}$ be the predictor that on input $\vz$, with oracle access
to $H$, reads all of $H$, computes $\pi^{i}_{H,z_i}$ --- it may, the sources being
fixed and predictors arbitrary --- and outputs a maximiser. Then
$\Prob{\mathsf{P}^{H}(\vz) = x_i} = \Exp_{(f,\zeta)}[m_i(f,\zeta)]$, which
unpredictability bounds by $\delta$. (iii) follows from $m_2 \le 1$ and (ii). For
(iv), $m_1m_2 > \vartheta$ with both factors in $[0,1]$ forces
$m_i > \sqrt{\vartheta}$ for some $i$; Markov and (ii) give
$\Prob{m_i > \sqrt{\vartheta}} \le \delta/\sqrt{\vartheta}$.
\end{proof}

Fact (i) is where the second source is spent: a single cell of $[N]^{2}$ carries
mass $m_1m_2$, a product of two independently small numbers, whereas one
unpredictable source over $[N]^{2}$ gives only $\delta$ per cell.

\begin{remark}[{$\Exp[m_1m_2] \le \delta^{2}$ is false}]\label{rem:m1m2}
This is needed later to justify an exponent. Let $2 \le M \le N$, $\delta := 1/M$,
let $E$ be the event $f(1,1) = 1$, of probability $\delta$, and let both sources
output $x_i := 1$ if $E$ holds and $x_i$ uniform on $[N]$ otherwise, leaking the
empty string. The coins are independent, so the pair is split, and
$\Exp[m_i] = \delta + (1-\delta)/N \le 2\delta$, so the pair is
$2\delta$-unpredictable; yet $m_1m_2 = 1$ on $E$, so $\Exp[m_1m_2] \ge \delta$.
Splitness constrains the coins, not the two sources' common dependence on $f$.
\end{remark}

\begin{lemma}[flattening]\label{lem:2}
Let $\pi$ be a distribution on $[N]$ with $\|\pi\|_\infty = m$ and
$k := \lfloor 1/m \rfloor$. Then $1 \le k \le N$, $1/k \le 2m$, and $\pi$ is a
convex combination of distributions uniform on subsets of size exactly $k$.
Consequently, for every $G : [N]^{2} \to [-1,1]$ and every $\pi^{1},\pi^{2}$ with
$\|\pi^{i}\|_\infty = m_i$, $k_i = \lfloor 1/m_i \rfloor$,
\[
  \bigl|\Exp_{\pi^{1} \otimes \pi^{2}}[G]\bigr| \le
  \max\bigl\{\, \bigl|\Exp_{\Unif(B_1) \otimes \Unif(B_2)}[G]\bigr|
  \ :\ |B_i| = k_i \,\bigr\}.
\]
\end{lemma}

\begin{proof}
$m \ge 1/N$ gives $k \le N$; $m \le 1$ gives $k \ge 1$. For $x \ge 1$,
$\lfloor x \rfloor \ge x/2$ --- for $1 \le x < 2$ because
$\lfloor x \rfloor = 1 \ge x/2$, for $x \ge 2$ because
$\lfloor x \rfloor > x-1 \ge x/2$ --- so with $x = 1/m$, $1/k \le 2m$. If $k = 1$
the decomposition is the identity $\pi = \sum_u \pi(u)\delta_u$ and no external
result is needed; assume $k \ge 2$. Since
$\|\pi\|_\infty \le 1/k = 2^{-\log k}$ with $\log k > 0$ and
$2^{\log k} = k \in \mathbb{N}$, Assumption~\ref{ass:flat} applies. The displayed
inequality follows because
$(\rho^1,\rho^2) \mapsto \Exp_{\rho^1 \otimes \rho^2}[G]$ is bilinear, so
substituting both decompositions writes the left side as a convex combination of
the quantities on the right.
\end{proof}

Note $G$ is arbitrary; it will carry an indicator
$\ind[\vu \notin \mathcal{S}]$ that does not factor over the coordinates, which is
harmless because the linear structure used is in $(\pi^1,\pi^2)$, not in $G$.

\begin{lemma}[second-moment sharpening]\label{lem:A}
With $m_1,m_2$ as in Lemma~\ref{lem:1}, and $t_c$ as in \eqref{eq:tc} below:
(a) $\Exp[\sqrt{m_1m_2}] \le \delta$;
(b) $\Prob{m_1m_2 > \vartheta} \le \delta/\sqrt{\vartheta}$ for
$\vartheta \in (0,1]$;
(c) pointwise, for $k_i = \lfloor 1/m_i \rfloor$ and every $c > 0$,
$t_c(k_1,k_2) \le \sqrt{L(m_1+m_2)} + \sqrt{2c}\sqrt{m_1m_2}$;
(d) $\Exp[t_c(k_1,k_2)] \le \sqrt{2L\delta} + \delta\sqrt{2c}$.
\end{lemma}

\begin{proof}
(a) $m_1,m_2$ are nonnegative and bounded by $1$, so all expectations are finite,
and by Cauchy--Schwarz and Lemma~\ref{lem:1}(ii) applied to each coordinate
separately,
$\Exp[\sqrt{m_1}\sqrt{m_2}] \le \sqrt{\Exp[m_1]}\sqrt{\Exp[m_2]} \le \delta$.
(b) $m_1m_2 > \vartheta$ iff $\sqrt{m_1m_2} > \sqrt{\vartheta}$; Markov and (a).
(c) $t_c^{2} = \frac{\ln(eN/k_1)}{2k_2} + \frac{\ln(eN/k_2)}{2k_1}
+ \frac{c}{2k_1k_2}$. Each $k_i \ge 1$, so $\ln(eN/k_i) \le L$; and
$1/k_i \le 2m_i$, so $\frac{1}{2k_i} \le m_i$ and
$\frac{1}{2k_1k_2} \le 2m_1m_2$. Hence
$t_c^{2} \le L(m_1+m_2) + 2c\,m_1m_2$, and
$\sqrt{a+b} \le \sqrt{a} + \sqrt{b}$ gives the claim.
(d) Take expectations in (c). By concavity of $\sqrt{\cdot}$ and
Lemma~\ref{lem:1}(ii),
$\Exp[\sqrt{L(m_1+m_2)}] \le \sqrt{L(\Exp m_1 + \Exp m_2)} \le \sqrt{2L\delta}$;
and $\sqrt{2c}\,\Exp[\sqrt{m_1m_2}] \le \delta\sqrt{2c}$ by (a).
\end{proof}

Part (a) is what buys the sharpening. Bounding
$\sqrt{L(\Exp m_1 + \Exp m_2) + 2c\,\Exp[m_1m_2]} \le \sqrt{2\delta(L+c)}$ by
concavity on the whole expression, as the cruder route does, gives
$\sqrt{2c\delta}$ where splitting first and applying Cauchy--Schwarz to the cross
term gives $\delta\sqrt{2c}$. The stronger hypothesis
$\Exp[m_1m_2] \le \delta^{2}$, which would give the same conclusion directly, is
false by Remark~\ref{rem:m1m2}, and (a) does not follow from
$\Exp[m_1m_2] \le \delta$ by Jensen, which runs the other way.

\section{Rectangle discrepancy}\label{sec:rectangle}

Put $L := \ln(eN)$ and, for $c > 0$,
\begin{equation}\label{eq:tc}
  t_c(k_1,k_2) := \sqrt{\frac{k_1 \ln\frac{eN}{k_1}
  + k_2 \ln\frac{eN}{k_2} + c}{2k_1k_2}} .
\end{equation}

A \emph{revealing rule of budget $s$} is a family
$(\mathcal{A}_\zeta)_{\zeta \in \mathcal{Z}}$ of procedures, where
$\mathcal{A}_\zeta$ inspects coordinates of $f$ one at a time --- the next a
function of $\zeta$ and the values already seen --- halts having inspected a set
$\mathcal{S}(f,\zeta)$ with $|\mathcal{S}(f,\zeta)| \le s$, and outputs a test
$\theta_{\zeta,f} : [M] \to \bits$ that is a function of $\zeta$ and the inspected
values.

The next two lemmas bound the same rectangle sum in two different ways, and the
difference between them is the whole difference between the two extraction bounds
of \S\ref{sec:extraction}. Lemma~\ref{lem:3} restricts the test class to the
$|\mathcal{Z}| < 2^{\sigma+2}$ tests a revealing rule can reach, and pays for the
revealed set. Lemma~\ref{lem:B} quantifies over all $2^{M}$ tests, and pays $M$
instead --- but owes nothing for revelation.

\begin{lemma}[restricted test class]\label{lem:3}
Let $N \ge 2$, $\gamma_0 \in (0,1)$ and $(\mathcal{A}_\zeta)$ a revealing rule. Put
$C_0 := (\sigma+2)\ln 2 + \ln \frac{4N^{2}}{\gamma_0}$. With probability at least
$1 - \gamma_0$ over $f \gets_{\$} \Fun$, simultaneously for all
$\zeta \in \mathcal{Z}$, all $k_1,k_2 \in [N]$ and all $B_i \subseteq [N]$ with
$|B_i| = k_i$, writing $R = B_1 \times B_2$,
\[
  \Bigl|\ \tfrac{1}{k_1k_2}\!\!\sum_{\vu \in R \setminus \mathcal{S}(f,\zeta)}\!\!
  \bigl(\theta_{\zeta,f}(f(\vu)) - p_{\theta_{\zeta,f}}\bigr)\Bigr|
  \ \le\ t_{C_0}(k_1,k_2).
\]
Moreover $\Exp[t_{C_0}] \le \sqrt{2\delta(L+C_0)}
\le \sqrt{8\delta(\sigma' + \log\gamma_0^{-1})}$.
\end{lemma}

\begin{proof}
\textbf{(1)} Fix $\zeta$. For $S_0 \subseteq [N]^{2}$ and $w \in [M]^{S_0}$ the
event $\mathcal{T}_{S_0,w} = \{\mathcal{S}(f,\zeta) = S_0 \wedge f|_{S_0} = w\}$ is
measurable with respect to $f|_{S_0}$: replaying $\mathcal{A}_\zeta$ on $w$
reproduces its whole inspection sequence. The coordinates of $f$ being i.i.d.\
uniform, conditioning on such an event leaves
$f|_{[N]^{2} \setminus S_0}$ i.i.d.\ uniform; and $\theta_{\zeta,f}$,
$p_{\theta_{\zeta,f}}$ are constant on $\mathcal{T}_{S_0,w}$.

\textbf{(2)} Fix also $k_1,k_2,B_1,B_2$ and condition on $\mathcal{T}_{S_0,w}$.
Write $\theta, p_\theta$ for the constants and
$n_0 := |R \setminus S_0| \le k_1k_2$. The variables
$\theta(f(\vu)) - p_\theta$, $\vu \in R \setminus S_0$, are independent, lie in an
interval of length $1$, and have mean $0$, since $f(\vu)$ is uniform on $[M]$ and
$\Prob{\theta(f(\vu)) = 1} = p_\theta$ exactly. Assumption~\ref{ass:binom} with
$\lambda := k_1k_2 t_{C_0}$ gives conditional failure probability at most
$2\exp(-2\lambda^{2}/n_0) \le 2\exp(-2t_{C_0}^{2}k_1k_2)$, uniformly in
$(S_0,w)$, hence also unconditionally.

\textbf{(3)} By \eqref{eq:tc},
$2\exp(-2t_{C_0}^{2}k_1k_2)
= 2 e^{-k_1 \ln\frac{eN}{k_1}} e^{-k_2 \ln\frac{eN}{k_2}} 2^{-(\sigma+2)}
\frac{\gamma_0}{4N^{2}}$. There are
$|\mathcal{Z}|\binom{N}{k_1}\binom{N}{k_2}$ tuples with the given $(k_1,k_2)$;
$\binom{N}{k} \le (eN/k)^{k}$ and $|\mathcal{Z}| < 2^{\sigma+2}$ cancel the three
exponential factors, leaving at most $\gamma_0/(2N^{2})$ per pair and $\gamma_0/2$
in total.

\textbf{(4)} $t_{C_0}^{2} \le \frac{L}{2k_2} + \frac{L}{2k_1}
+ \frac{C_0}{2k_1k_2}$ since $k_i \ge 1$, and $1/k_i \le 2m_i$ by
Lemma~\ref{lem:2} bounds the three terms by $Lm_2$, $Lm_1$, $2C_0m_1m_2$.

\textbf{(5)} By concavity of $\sqrt{\cdot}$ and Lemma~\ref{lem:1}(ii),(iii),
$\Exp[t_{C_0}] \le \sqrt{L(\Exp m_1 + \Exp m_2) + 2C_0\Exp[m_1m_2]}
\le \sqrt{2\delta(L+C_0)}$. Numerically
$L + C_0 = 1 + \ln N + (\sigma+2)\ln 2 + \ln 4 + 2\ln N + \ln\gamma_0^{-1}
\le 3\ln N + 0.694\sigma + 3.78 + \ln\gamma_0^{-1}$; as
$\sigma' \ge 2.885 \ln N$ and $\sigma' \ge 2$ for $N \ge 2$, this is at most
$3.63\sigma' + \ln\gamma_0^{-1}$, whence the claim.
\end{proof}

\begin{remark}[$M$ occurs nowhere in $t_{C_0}$]
The alphabet enters only through $p_\theta$, which cancels, and through the
i.i.d.\ structure of $\{f(\vu)\}$, which is alphabet-blind. There are $2^{M}$
possible tests, but once the observer is fixed and the revealed values conditioned
away, only $|\mathcal{Z}| < 2^{\sigma+2}$ of them are reached.
\end{remark}

\begin{lemma}[universal test class]\label{lem:B}
Let $N,M \ge 2$, $\gamma_0 \in (0,1)$, and put
$C_1 := M \ln 2 + \ln \frac{4N^{2}}{\gamma_0}$. With probability at least
$1-\gamma_0$ over $f \gets_{\$} \Fun$, simultaneously for all $k_1,k_2 \in [N]$,
all $B_i$ with $|B_i| = k_i$, and \emph{all} $\theta : [M] \to \bits$, writing
$R := B_1 \times B_2$,
\[
  \Bigl|\frac{1}{k_1k_2}\sum_{\vu \in R}
  \bigl(\theta(f(\vu)) - p_\theta\bigr)\Bigr| \le t_{C_1}(k_1,k_2).
\]
\end{lemma}

\begin{proof}
\textbf{(1)} Fix $k_1,k_2$, sets $B_i$ with $|B_i| = k_i$, and a function $\theta$
--- all fixed before $f$ is drawn. For $\vu \in R$ put
$Y_{\vu} := \theta(f(\vu)) - p_\theta$. The coordinates of $f$ are i.i.d.\ uniform
on $[M]$, so the $k_1k_2$ variables $Y_{\vu}$ are independent; each has mean zero
exactly, since $\Prob{\theta(f(\vu)) = 1} = |\theta^{-1}(1)|/M = p_\theta$; and each
lies in $\{-p_\theta, 1-p_\theta\}$, an interval of length $1$.
Assumption~\ref{ass:binom} with $\lambda := k_1k_2 t_{C_1}$ gives failure
probability at most $2\exp(-2t_{C_1}^{2}k_1k_2)$.

\textbf{(2)} By \eqref{eq:tc},
$2t_{C_1}^{2}k_1k_2 = k_1 \ln\frac{eN}{k_1} + k_2 \ln\frac{eN}{k_2} + C_1$, so that
bound equals
$2(\frac{eN}{k_1})^{-k_1}(\frac{eN}{k_2})^{-k_2}2^{-M}\frac{\gamma_0}{4N^{2}}$.

\textbf{(3)} For the fixed pair $(k_1,k_2)$ the number of tuples
$(B_1,B_2,\theta)$ is
$\binom{N}{k_1}\binom{N}{k_2}2^{M} \le (\frac{eN}{k_1})^{k_1}(\frac{eN}{k_2})^{k_2}2^{M}$,
by Assumption~\ref{ass:binom} and the fact that there are exactly $2^{M}$
functions $[M] \to \bits$. The three exponential factors cancel exactly, leaving at
most $\gamma_0/(2N^{2})$ per pair and $\gamma_0/2 \le \gamma_0$ over the at most
$N^{2}$ pairs.
\end{proof}

\begin{remark}[the quantifier order is the mechanism]\label{rem:qorder-B}
The good event is a property of $f$ alone, and on it the bound holds for every
$\theta$ fixed in advance; substituting the particular $\theta_{\zeta,f}$, which
does depend on $f$, is legitimate because a statement ``for all $\theta$,
$\Phi(f,\theta) \le t$'' may be instantiated at any $\theta$ whatsoever. This is
why no conditioning on inspected values is needed, why the sum runs over all of
$R$, and why no revealing rule, budget, excised set or term
$\Pi(\mathcal{S})$ appears. The trade against Lemma~\ref{lem:3} is
$\sigma+2 \rightsquigarrow M$ inside a budget paid at rate $\delta\sqrt{2C_1}$,
i.e.\ $\delta\sqrt{\sigma} \rightsquigarrow \delta\sqrt{M}$ in the final bound,
and \emph{not} $\sqrt{\sigma\delta} \rightsquigarrow \sqrt{M\delta}$.
\end{remark}

\begin{remark}[where one source would break both]\label{rem:onesource}
With a single source emitting a point of $[N]^{2}$, Lemma~\ref{lem:1}(i) is
unavailable and Lemma~\ref{lem:2} flattens onto an arbitrary subset
$B \subseteq [N]^{2}$ of size $k = \lfloor 1/m \rfloor$, not onto a rectangle. Step
(3) would then range over $\binom{N^{2}}{k} \le (eN^{2}/k)^{k}$ sets, forcing
$t^{2} \ge \frac12 \ln(eN^{2}/k) \ge \frac12$ since $k \le N^{2}$, so
$t \ge 1/\sqrt2$ and the conclusion is vacuous, independently of $\delta$, $M$,
$\sigma$. Lemma~\ref{lem:A}(a) collapses likewise, its Cauchy--Schwarz step being
applied across the two coordinates. The failure occurs before any bound on the
probability that a fixed set contains $\vx$ is invoked, which is what the
conjecture's own single-source counterexample requires.
\end{remark}

\section{The observer as a revealing rule}\label{sec:revealing}

\begin{definition}[challenge resolution]\label{def:res}
$D$ has \emph{challenge resolution $M'$} if there is $h : [M] \to [M']$ such that
for \emph{every coin string $\rho$ of $D$}, every $f$, every $\zeta$, and all
$v,v'$ with $h(v) = h(v')$, the runs $D^{f}_\rho(v,\zeta)$ and
$D^{f}_\rho(v',\zeta)$ issue the same sequence of query \emph{positions}; the output
bit is unconstrained. Every $D$ has resolution $M$; resolution $1$ means that, for
each fixed coin string, the query positions ignore the challenge value. We may take
$h$ surjective without loss.
\end{definition}

The coin quantifier is part of the definition and not a technicality. The weaker
condition that the query positions be challenge-independent \emph{in law} defines a
strictly larger class: at $N = M = 2$, $q = 1$, the observer that flips $\rho$ and
queries $(1,1)$ when $v \oplus \rho = 1$ and $(2,2)$ otherwise has query position
uniform on $\{(1,1),(2,2)\}$ for both challenge values, yet each of its two coin
fixings has resolution $2$.

\begin{lemma}\label{lem:4}
Let $D$ be deterministic, $q$-query, challenge-oblivious, of resolution $M'$. Then
the procedures ``$\mathcal{A}_\zeta$: for $j = 1,\dots,M'$ run $D^{f}(v_j,\zeta)$
for a fixed $v_j \in h^{-1}(j)$, inspecting each queried cell not already known''
form a revealing rule of budget $s = \min(qM',N^{2})$, with
$\theta_{\zeta,f}(v) := D^{f}(v,\zeta)$.
\end{lemma}

\begin{proof}
$\mathcal{A}_\zeta$ inspects one coordinate at a time, the next a function of
$\zeta$ and the values already seen: the loop and the representatives are fixed in
advance, and inside run $j$ the next position is a function of $(v_j,\zeta)$ and
previous answers. Each of the $M'$ runs issues at most $q$ positions, so at most
$qM'$ cells, and at most $N^{2}$ exist. Finally, for arbitrary $v$ with
$j = h(v)$, the run $D^{f}(v,\zeta)$ issues exactly the positions run $j$ issued,
all inspected and recorded, so its whole execution --- hence its output bit --- is
reconstructible from $(v,\zeta,\mathcal{S},f|_{\mathcal{S}})$. Therefore the entire
function $\theta_{\zeta,f}$, and $p_{\theta_{\zeta,f}}$, is determined by
$(\zeta,\mathcal{S},f|_{\mathcal{S}})$.
\end{proof}

\begin{lemma}[mass of the revealed set]\label{lem:C}
Let $\mathcal{S}(f,\zeta) \subseteq [N]^{2}$ be a function of $(f,\zeta)$ with
$|\mathcal{S}| \le s$. Then for $s \ge 1$
\[
  \Exp\bigl[\min(\Pi_{f,\zeta}(\mathcal{S}(f,\zeta)),1)\bigr] \le
  \mu'(s) := \min\bigl(s\delta,\ 2(s\delta^{2})^{1/3},\ 1\bigr),
\]
and the left side is $0$ for $s = 0$.
\end{lemma}

\begin{proof}
By Lemma~\ref{lem:1}(i),
$\Pi(\mathcal{S}) = \sum_{\vu \in \mathcal{S}} \pi^1(u_1)\pi^2(u_2) \le s\,m_1m_2$;
expectations and Lemma~\ref{lem:1}(iii) give the arm $s\delta$, and the arm $1$ is
trivial. For the third, fix $\vartheta \in (0,1]$, bound by $s\vartheta$ on
$\{m_1m_2 \le \vartheta\}$ and by $1$ off it, and use Lemma~\ref{lem:A}(b):
$\Exp[\min(\Pi(\mathcal{S}),1)] \le g(\vartheta) := s\vartheta + \delta\vartheta^{-1/2}$.
Then $g'(\vartheta) = s - \frac12 \delta\vartheta^{-3/2}$ vanishes at
$\vartheta^{*} = (\delta/2s)^{2/3}$, which lies in $(0,1]$ since
$\delta \le 1 \le 2s$, and $g'' > 0$; and
$g(\vartheta^{*}) = (2^{-2/3} + 2^{1/3})(s\delta^{2})^{1/3}
\le 1.89(s\delta^{2})^{1/3} \le 2(s\delta^{2})^{1/3}$.
\end{proof}

The arm $2(s\delta^{2})^{1/3}$ cannot be improved to $s\delta^{2}$, since
$\Exp[m_1m_2] \le \delta^{2}$ is false (Remark~\ref{rem:m1m2}). The exponent
$1/3$ is exactly what the $\sqrt{\vartheta}$ of Lemma~\ref{lem:A}(b) produces, and
the bound is tight at $s = 1$ on that counterexample.

\section{Bounds on the extraction advantage}\label{sec:extraction}

\begin{lemma}\label{claim:61}
For deterministic $D$ with associated $\theta_{\zeta,f}$,
$\Prob{\Real = 1} = \Exp_{(f,\zeta)}\bigl[\Exp_{\vx \sim \Pi_{f,\zeta}}
[\theta_{\zeta,f}(f(\vx))]\bigr]$ and
$\Prob{\Real_0 = 1} = \Exp_{(f,\zeta)}[p_{\theta_{\zeta,f}}]$.
\end{lemma}

\begin{proof}
Condition on $(H,\vz) = (f,\zeta)$. In both experiments $D$'s oracle is $f$, so its
output on challenge $v$ is $\theta_{\zeta,f}(v)$. In $\Real$ the challenge is
$f(\vx)$ with $\vx \sim \Pi_{f,\zeta}$ by Lemma~\ref{lem:1}; in $\Real_0$ it is
uniform and independent, with $\Exp_{v \gets_{\$} [M]}[\theta(v)] = p_\theta$.
\end{proof}

\begin{proposition}\label{prop:62}
For every $\gamma_0 \in (0,1)$ and deterministic $q$-query challenge-oblivious $D$
of resolution $M'$,
\[
  \bigl|\Prob{\Real = 1} - \Prob{\Real_0 = 1}\bigr| \le
  \gamma_0 + \sqrt{8\delta(\sigma' + \log\gamma_0^{-1})}
  + \mu'\bigl(\min(qM',N^{2})\bigr).
\]
\end{proposition}

\begin{proof}
By Lemma~\ref{claim:61} and the triangle inequality the left side is at most
$\Exp[\Delta]$ with
$\Delta(f,\zeta) := |\Exp_{\Pi}[\theta(f(\vx)) - p_\theta]|$. Put
$G_{f,\zeta}(\vu) := (\theta(f(\vu)) - p_\theta)\ind[\vu \notin \mathcal{S}(f,\zeta)]
\in [-1,1]$, so that
$\Delta \le |\Exp_{\Pi}[G_{f,\zeta}]| + \Pi(\mathcal{S})$. Lemma~\ref{lem:2} bounds
$|\Exp_{\Pi}[G]|$ by the largest
$|\Exp_{\Unif(B_1) \otimes \Unif(B_2)}[G]|$ over $|B_i| = k_i = \lfloor 1/m_i\rfloor$,
and for such a rectangle that quantity equals
$|\frac{1}{k_1k_2}\sum_{\vu \in R \setminus \mathcal{S}}(\theta(f(\vu)) - p_\theta)|$,
which on the good event of Lemma~\ref{lem:3} is at most $t_{C_0}(k_1,k_2)$. As
$\Delta \le 1$ always,
$\Exp[\Delta] \le \Prob{\mathcal{E}^{c}} + \Exp[t_{C_0}]
+ \Exp[\min(\Pi(\mathcal{S}),1)]$, and the three terms are bounded by $\gamma_0$,
by Lemma~\ref{lem:3}, and by $\mu'(s)$ via Lemmas~\ref{lem:4} and~\ref{lem:C}.
\end{proof}

\begin{theorem}[$M$-free, resolution-restricted]\label{thm:A}
$\kna(q) \le 5\sqrt{\sigma'\delta} + q\delta$, and
$\kappa(0) \le 5\sqrt{\sigma'\delta}$.
\end{theorem}

\begin{proof}
Both suprema are over observers permitted to randomise, whereas
Proposition~\ref{prop:62} is stated for deterministic $D$; by Lemma~\ref{lem:0} the
suprema are attained on deterministic observers, and $D'$ inherits $D$'s
resolution, so restricting to deterministic resolution-$1$ observers loses nothing.
If $\delta = 1$ then $\sqrt{\sigma'\delta} \ge \sqrt2 > 1$ and there is nothing to
prove; else take $\gamma_0 := \delta$. Since $\delta \ge 1/N$ we get
$\log\gamma_0^{-1} \le \log N \le \sigma'/2$, so the middle term is at most
$\sqrt{12\sigma'\delta} \le 3.47\sqrt{\sigma'\delta}$, and
$\gamma_0 = \delta \le \sqrt{\sigma'\delta}$ as $\delta \le 1 \le \sigma'$.
Resolution $1$ gives $s \le q$ and $\mu'(q) \le q\delta$. For $q = 0$, $s = 0$ and
every observer has resolution $1$ vacuously.
\end{proof}

\begin{theorem}[$q$-free, every observer]\label{thm:D}
Let $N,M \ge 2$ and $\gamma_0 \in (0,1)$. For \emph{every} challenge-oblivious
observer $D$ --- of any query count, bounded or unbounded, and any challenge
resolution ---
\[
  \bigl|\Prob{\Real = 1} - \Prob{\Real_0 = 1}\bigr| \le
  \gamma_0 + \sqrt{2\delta \ln(eN)}
  + \delta\sqrt{2M\ln 2 + 2\ln \tfrac{4N^{2}}{\gamma_0}} .
\]
\end{theorem}

\begin{proof}
By Lemma~\ref{lem:0} we may assume $D$ deterministic; Lemma~\ref{lem:0} imposes no
bound on the query count and nothing below refers to one. With $\theta_{\zeta,f}$
and $p_\theta$ as in Lemma~\ref{claim:61}, set
$\Delta(f,\zeta) := |\Exp_{\vu \sim \Pi_{f,\zeta}}
[\theta_{\zeta,f}(f(\vu)) - p_{\theta_{\zeta,f}}]|$; by Lemma~\ref{claim:61} and the
triangle inequality the left side is at most $\Exp[\Delta]$.

Fix $(f,\zeta)$ in the support and put
$G_{f,\zeta}(\vu) := \theta_{\zeta,f}(f(\vu)) - p_{\theta_{\zeta,f}} \in [-1,1]$. By
Lemma~\ref{lem:1} the law $\Pi_{f,\zeta}$ is a product, so Lemma~\ref{lem:2}
applies with this $G$ and gives
\[
  \Delta(f,\zeta) \le \max_{|B_i| = k_i}
  \Bigl|\frac{1}{k_1k_2}\sum_{\vu \in B_1 \times B_2}
  \bigl(\theta_{\zeta,f}(f(\vu)) - p_{\theta_{\zeta,f}}\bigr)\Bigr| .
\]
Let $\mathcal{E}$ be the good event of Lemma~\ref{lem:B}, a property of $f$ alone,
of probability at least $1 - \gamma_0$. On $\mathcal{E}$ the maximum is at most
$t_{C_1}(k_1,k_2)$: the pair $(k_1,k_2)$ lies in $[N] \times [N]$ by
Lemma~\ref{lem:2}, the maximising $B_i$ are among the sets quantified over, and
$\theta_{\zeta,f}$ is one particular element of $\bits^{[M]}$, over all of which
Lemma~\ref{lem:B} quantifies --- this is Remark~\ref{rem:qorder-B}. Since
$G \in [-1,1]$ gives $\Delta \le 1$ pointwise, and $t_{C_1} \ge 0$,
\[
  \Exp[\Delta] \le \Prob{\mathcal{E}^{c}} + \Exp[\ind_{\mathcal{E}} t_{C_1}]
  \le \gamma_0 + \Exp[t_{C_1}] \le \gamma_0 + \sqrt{2L\delta} + \delta\sqrt{2C_1},
\]
by Lemma~\ref{lem:A}(d) at $c = C_1$.
\end{proof}

For $M = 1$ the set $\Fun$ is a singleton and $\Real, \Real_0$ are the same
experiment, so the advantage is $0$ and the excluded case needs no bound.

\begin{corollary}\label{cor:Dp}
For every $q \in \mathbb{N} \cup \{0\}$,
$\kappa(q) \le 5\sqrt{\sigma'\delta} + 2\delta\sqrt{M}$, with no restriction on the
observer.
\end{corollary}

\begin{proof}
If $\delta = 1$ then $5\sqrt{\sigma'\delta} \ge 5\sqrt2 > 1 \ge \kappa(q)$; assume
$\delta < 1$ and take $\gamma_0 := \delta$. Recall $\delta \ge 1/N$, $N \ge 2$, and
$\sigma' \ge 2\log N \ge 2$, so $\ln N \le \frac{\ln 2}{2}\sigma' \le 0.3466\sigma'$
and $1 \le \sigma'/2$. Then $2\delta L = 2\delta(1 + \ln N) \le 1.6932\sigma'\delta$,
so $\sqrt{2\delta L} \le 1.3012\sqrt{\sigma'\delta}$; and, using
$\ln(1/\delta) \le \ln N$,
\[
  2C_1 = 2M\ln 2 + 2\ln 4 + 4\ln N + 2\ln\tfrac1\delta
  \le 1.3863M + 2.7726 + 6\ln N \le 1.3863M + 3.4657\sigma',
\]
whence
$\delta\sqrt{2C_1} \le 1.1774\delta\sqrt{M} + 1.8617\delta\sqrt{\sigma'}
\le 1.1774\delta\sqrt{M} + 1.8617\sqrt{\sigma'\delta}$, using
$\delta\sqrt{\sigma'} = \sqrt{\delta}\sqrt{\sigma'\delta} \le \sqrt{\sigma'\delta}$.
Finally $\gamma_0 = \delta \le \sqrt{\sigma'\delta}$ as $\delta \le 1 \le \sigma'$.
Adding, $(1 + 1.3012 + 1.8617)\sqrt{\sigma'\delta} + 1.1774\delta\sqrt{M}
\le 5\sqrt{\sigma'\delta} + 2\delta\sqrt{M}$.
\end{proof}

\begin{theorem}[the revealing-rule arm, collected]\label{thm:ApBp}
For every $\gamma_0$ and every deterministic $q$-query challenge-oblivious $D$ of
resolution $M'$, Proposition~\ref{prop:62} with Lemma~\ref{lem:A}(d) in place of
Lemma~\ref{lem:3}'s cruder expectation bound gives
\[
  \kappa(q) \le \sqrt{2\delta\ln(eN)} + 4\delta\sqrt{\sigma'}
  + \mu'\bigl(\min(qM,N^{2})\bigr),
\]
and the same with $\mu'(q)$ at $M' = 1$, which bounds $\kna(q)$.
\end{theorem}

\begin{proof}
Run Proposition~\ref{prop:62} with $\Exp[t_{C_0}]$ bounded by
Lemma~\ref{lem:A}(d) at $c = C_0$ rather than by concavity on the whole
expression, and take $\gamma_0 = \delta$. Then $\ln(1/\delta) \le \ln N$ and
$2C_0 \le 1.3863\sigma + 5.5452 + 6\ln N \le 6.238\sigma'$, giving
$\delta\sqrt{2C_0} \le 2.498\delta\sqrt{\sigma'}$; and
$\gamma_0 = \delta \le 0.708\delta\sqrt{\sigma'}$, the two together at most
$3.21\delta\sqrt{\sigma'} \le 4\delta\sqrt{\sigma'}$.
\end{proof}

\begin{corollary}[the two arms]\label{cor:Dpp}
$\kappa(q) \le 5\sqrt{\sigma'\delta}
+ \min\{\mu'(\min(qM,N^{2})),\ 2\delta\sqrt{M}\}$, the first arm granting
Theorem~\ref{thm:ApBp} and the second unconditional by
Corollary~\ref{cor:Dp}. Both leading parts are at most $5\sqrt{\sigma'\delta}$:
$4.164$ for the second, and
$1.302\sqrt{\sigma'\delta} + 3.21\delta\sqrt{\sigma'} \le 4.512\sqrt{\sigma'\delta}$
for the first, using $\delta\sqrt{\sigma'} \le \sqrt{\sigma'\delta}$.
\end{corollary}

\begin{remark}[which arm wins]\label{rem:whicharm}
Suppose $q \ge 1$, $M \ge 4$, and $2\delta\sqrt{M} < 1$. Then the second arm is the
smaller: against $\mu'$'s arm $1$ immediately; against $2(s\delta^{2})^{1/3}$ with
$s \le qM$ because $2\delta\sqrt{M} \le 2(qM\delta^{2})^{1/3}$ iff
$\delta\sqrt{M} \le q$, and $\delta\sqrt{M} < 1/2 < q$; against
$s\delta \le qM\delta$ because $q\sqrt{M} \ge 2$. The cap $s = N^{2}$ never binds
usefully, since $\mu'(N^{2}) = 1$: both $N^{2}\delta \ge N \ge 2$ and
$2(N\delta)^{2/3} \ge 2$, as $\delta \ge 1/N$. So for every $q \ge 1$ and
$M \ge 4$ the second arm dominates wherever either says anything; the first
survives only at $q = 0$, where $\mu'(0) = 0$ is exact and $\delta\sqrt{M}$ is
pure loss, and in the corner $M \in \{2,3\}$.
\end{remark}

\section{A lower bound}\label{sec:lower}

\begin{proposition}\label{prop:F}
Let $N \ge 2$ and $2 \le M \le N^{2}/2$. Let $S_1, S_2$ each draw
$x_i \gets_{\$} [N]$ on private coins and output $(x_i,\varepsilon)$. Then the pair
is split, $\sigma_1 = \sigma_2 = 0$, $\sigma' = 2\log N$, and it is
$\delta$-unpredictable with $\delta = 1/N$. There is a deterministic
challenge-oblivious $D^{*}$, making $N^{2}$ queries and of challenge resolution
$1$, with
$|\Prob{\Real = 1} - \Prob{\Real_0 = 1}| \ge \delta\sqrt{M}/(4\sqrt2)$.
\end{proposition}

\begin{proof}
\textbf{(1)} The coins are independent, so the pair is split. Here
$\vz = (\varepsilon,\varepsilon)$ always and each $x_i$ is uniform on $[N]$ and
independent of $H$ and of $x_{3-i}$. A predictor is an arbitrary function of
$(\vz,H)$, so its output is independent of $x_i$ and
$\Prob{\mathsf{P}^{H}(\vz) = x_i} = 1/N$ for every $\mathsf{P}$.

\textbf{(2)} $D^{*}$ queries all $N^{2}$ cells, so holds $f$; forms
$\nu(v) := |f^{-1}(v)|/N^{2}$; and on challenge $v$ outputs
$\ind[\nu(v) > 1/M]$. Its query positions are the same for every challenge value,
so its resolution is $1$.

\textbf{(3)} Condition on $H = f$. In $\Real$, $\vx$ is uniform on $[N]^{2}$ and
independent of $f$, so the challenge $f(\vx)$ has law exactly $\nu$; in $\Real_0$
it is uniform on $[M]$. With $A_f := \{v : \nu(v) > 1/M\}$,
\[
  \Prob{\Real = 1 \mid f} - \Prob{\Real_0 = 1 \mid f}
  = \sum_v \bigl(\nu(v) - \tfrac1M\bigr)^{+}
  = \SD\bigl(\nu, \Unif([M])\bigr) \ \ge 0,
\]
so no cancellation occurs on averaging and the advantage equals
$\Exp_f[\SD(\nu,\Unif([M]))]$.

\textbf{(4)} \emph{Claim: let $Y = \sum_{i=1}^{n} Y_i$ with $Y_i$ i.i.d.,
$\Exp Y_1 = 0$, $|Y_1| \le 1$, and $s^{2} := \Exp Y^{2} = n\Exp Y_1^{2}$. If
$s \ge 1$ then $\Exp|Y| \ge s/2$.} Discarding the terms with an index to an odd
power, which vanish,
$\Exp Y^{4} = n\Exp Y_1^{4} + 3n(n-1)(\Exp Y_1^{2})^{2}
\le n\Exp Y_1^{2} + 3n^{2}(\Exp Y_1^{2})^{2} = s^{2} + 3s^{4}$, using
$\Exp Y_1^{4} \le \Exp Y_1^{2}$, valid as $|Y_1| \le 1$. H\"older with exponents
$3/2$ and $3$ gives
$s^{2} = \Exp[|Y|^{2/3}|Y|^{4/3}] \le (\Exp|Y|)^{2/3}(\Exp Y^{4})^{1/3}$, so
$\Exp|Y| \ge s^{3}/\sqrt{s^{2}+3s^{4}} \ge s^{3}/(2s^{2}) = s/2$ for $s \ge 1$.

\textbf{(5)} Put $n := N^{2}$, $p := 1/M$, $X_v := |f^{-1}(v)|$. Over uniform $f$
the summands $\ind[f(\vu) = v]$ are i.i.d.\ Bernoulli$(p)$, so
$X_v \sim \mathrm{Bin}(n,p)$ and each centred summand lies in $[-p,1-p]$, of
absolute value at most $1$. As
$\SD(\nu,\Unif([M])) = \frac{1}{2n}\sum_v |X_v - np|$ and the $X_v$ are
identically distributed, $\Exp_f[\SD] = \frac{M}{2n}\Exp|X - np|$. Here
$s^{2} = np(1-p) = \frac{N^{2}}{M}(1 - \frac1M) \ge \frac{N^{2}}{2M}$ for
$M \ge 2$, and $s^{2} \ge 1$ since $M \le N^{2}/2$, so (4) gives
$\Exp|X - np| \ge \frac{N}{2\sqrt{2M}}$ and
$\Exp_f[\SD] \ge \frac{M}{2N^{2}} \cdot \frac{N}{2\sqrt{2M}}
= \frac{\sqrt{M}}{4\sqrt2 N}$.
\end{proof}

So on the diagonal $\delta = 1/N$, Corollary~\ref{cor:Dp}'s $M$-dependence is
pinned to within $8\sqrt2 < 12$. The family fixes $\delta = 1/N$ and
$2 \le M \le N^{2}/2$, and licenses nothing off that diagonal --- nor anything at
$q < N^{2}$, which matters in \S\ref{sec:h2}.

\section{Presampling, in oracle-indexed consistent form}\label{sec:presampling}

\begin{lemma}\label{lem:P}
Let $N,M \ge 2$, $P \in \mathbb{N}$, $\gamma \in (0,1)$. There is a family
$\mathcal{Y}^{P,\gamma} = \{Y_{f,\zeta}\}$ indexed by $f \in \Fun$ and
$\zeta \in (\bits^{*})^{2}$ --- with $Y_{f,\zeta}$ uniform on $\Fun$ at every
$\zeta$ outside the support of $\vz$, so the family is total, no experiment
reaching such a $\zeta$ --- depending only on $(S_1,S_2,P,\gamma)$, in which every
$Y_{f,\zeta}$ is a $P$-mixture consistent with $f$ --- indeed a single
$P$-bit-fixing source --- such that for every $q$ and every $q$-query observer
taking $\vz$ and an independent uniform value of $[M]$,
\[
  \bigl|\Prob{\Real_0 = 1} - \Prob{\Dec_0 = 1}\bigr| \le
  \frac{q(\sigma + 2 + 2\log\gamma^{-1})}{P} + 2\gamma .
\]
\end{lemma}

\begin{proof}
\textbf{Step 1.} For $\zeta$ with $\Prob{\vz = \zeta} > 0$ let $\mu_\zeta$ be the
law of $H$ given $\vz = \zeta$ and $S_\zeta := N^{2}\log M - H_\infty(\mu_\zeta)$.
Since
$\mu_\zeta(f) = \Prob{H = f}\Prob{\vz = \zeta \mid H = f}/\Prob{\vz = \zeta}
\le (|\Fun|\Prob{\vz = \zeta})^{-1}$, we get
$S_\zeta \le \log(1/\Prob{\vz = \zeta})$. With
$\bar S := \sigma + 2 + \log\gamma^{-1}$ and
$\mathcal{B} := \{\zeta : S_\zeta > \bar S\}$, every $\zeta \in \mathcal{B}$ has
$\Prob{\vz = \zeta} < \gamma 2^{-(\sigma+2)}$, so
$\Prob{\vz \in \mathcal{B}} < |\mathcal{Z}|\gamma 2^{-(\sigma+2)} < \gamma$.

\textbf{Step 2.} Fix $\zeta$ and set
$\delta_\zeta := (S_\zeta + \log\gamma^{-1})/(P\log M)$. If $\delta_\zeta > 1$ the
asserted bound exceeds $q\log M \ge q$, hence is vacuous for $q \ge 1$, and for
$q = 0$ both experiments give the observer a uniform challenge and no oracle
access; set $Y_{f,\zeta}$ uniform. Otherwise apply Assumption~\ref{ass:cdgs2} to
$\mu_\zeta$ with parameter $\delta_\zeta$; since
$(S_\zeta + \log\gamma^{-1})/(\delta_\zeta \log M) = P$ it yields
$\mu_\zeta = \sum_j \lambda_j X_j + \lambda_{\mathrm{fin}}Y_{\mathrm{fin}}$ with
$\lambda_{\mathrm{fin}} \le \gamma$, each $X_j$ a $(P,1-\delta_\zeta)$-dense
source, all supports pairwise disjoint, and each $X_j$ fixing its set $I_j$,
$|I_j| \le P$, to values every $f$ in its support takes. Let $Y_j$ be the
corresponding $P$-bit-fixing source and define $Y_{f,\zeta} := Y_j$ for the unique
$j$ with $f \in \supp X_j$, and $Y_{f,\zeta} :=$ uniform otherwise. Consistency is
immediate; the construction reads only $(f,\zeta)$ and the decomposition, which
depends only on $(S_1,S_2,P,\gamma)$. Write $J = J(f,\zeta)$.

\textbf{Step 3.} Fix $\zeta \notin \mathcal{B}$ with $\delta_\zeta \le 1$ and
condition on $J = j$. Since $\lambda_j X_j$ is the restriction of $\mu_\zeta$ to
$\supp X_j$, the conditional law of $H$ is exactly $X_j$; in $\Dec_0$ the oracle is
drawn freshly from $Y_j$. The challenge is uniform and independent, so conditioning
on it leaves a fixed $q$-query adaptive distinguisher between $X_j$ and $Y_j$;
Assumption~\ref{ass:cdgs3} bounds its advantage by
$q\delta_\zeta \log M = q(S_\zeta + \log\gamma^{-1})/P$, and averaging over the
challenge preserves this.

\textbf{Step 4.} Averaging over $j$ and bounding the residue event, of conditional
probability at most $\gamma$, by $1$, then averaging over $\zeta$ and bounding
$\mathcal{B}$'s contribution by $\Prob{\vz \in \mathcal{B}} < \gamma$ while using
$S_\zeta \le \bar S$ off $\mathcal{B}$, gives the claim.
\end{proof}

For $M = 1$, $\Fun$ is a singleton, so $\Real_0$ and $\Dec_0$ coincide and the
bound holds with left side $0$; the construction divides by $\log M$ and is stated
for $M \ge 2$ only.

\section{From extraction to decomposition}\label{sec:decomposition}

\begin{theorem}\label{thm:C}
With $\mathcal{Y} := \mathcal{Y}^{P,\gamma/2}$, for every $q$ and every $q$-query
challenge-oblivious $D$, writing
$\varepsilon(D) := |\Prob{\Real = 1} - \Prob{\Real_0 = 1}|$ for that same $D$,
\[
  \Adv_{\mathcal{Y},D} \le \varepsilon(D)
  + \frac{2(\sigma' + \log\gamma^{-1})q}{P} + \gamma + P\delta + q\delta .
\]
\end{theorem}

\begin{proof}
Interpolate $\mathsf{G}_0 = \Real$, $\mathsf{G}_1 = \Real_0$,
$\mathsf{G}_2 = \Dec_0$, $\mathsf{G}_3 = \Dec$.

$|\mathsf{G}_0 - \mathsf{G}_1| = \varepsilon(D)$ by definition, for the fixed $D$
the theorem is instantiated at; no supremum is taken here.

$|\mathsf{G}_1 - \mathsf{G}_2|$: Lemma~\ref{lem:P} with slack $\gamma/2$ gives
$q(\sigma + 2 + 2\log(2/\gamma))/P + \gamma$; as $N \ge 2$ gives
$\sigma + 2 \le \sigma'$ and $\sigma' \ge 2$, the numerator is at most
$2\sigma' + 2\log\gamma^{-1}$. For $q \ge P$ that term already exceeds
$2\sigma' > 1$, so no hypothesis $q < P$ is needed.

$|\mathsf{G}_2 - \mathsf{G}_3| \le P\delta + q\delta$: draw $H$, $(\vx,\vz)$ and
$J$ with fixed set $I_J$; let $H^{\circ} \sim Y_J$ and $U \gets_{\$} [M]$ be
independent of everything else; let $H^{*}$ be $H^{\circ}$ with its value at $\vx$
overwritten by $U$ when $\vx \notin I_J$. As $J$ is a function of $(H,\vz)$,
$H^{\circ}$ is independent of $\vx$ given $(H,\vz,J)$; conditioning on
$\vx = \vu \notin I_J$, the source $Y_J$ is uniform and independent at $\vu$, so
both $H^{\circ}$ and $H^{*}$ are distributed as $Y_J$. Running $\mathsf{G}_3$ with
oracle $H^{*}$ and $\mathsf{G}_2$ with oracle $H^{\circ}$ and challenge $U$, the
two executions receive the same input on $\{\vx \notin I_J\}$, since there
$H^{*}(\vx) = U$, and their oracles agree off $\vx$. So they diverge only if
$\vx \in I_J$ or $D$ queries $\vx$; being identical up to the step before a first
query at $\vx$, that event has equal probability in both and may be measured in
$\mathsf{G}_2$. Lemma~\ref{lem:hit} gives $\Prob{\vx \in I_J} \le P\delta$ --- its
hypotheses hold, $\mathcal{Y}$ depending on $(S_1,S_2,P,\gamma)$ alone, being
indexed by $(f,\zeta)$, and having its component drawn independently of $\vx$ given
$(H,\vz)$, since $J$ is a function of $(H,\vz)$. Lemma~\ref{lem:query} gives
$\Prob{\vx \in Q} \le q\delta$ in $\mathsf{G}_2$, its predictor being simulable
because the challenge there is uniform and independent.
\end{proof}

\begin{remark}[where the $P$-cap enters, and that it is one step]\label{rem:onestep}
$P\delta$ enters the entire chain at exactly the step above, via
Lemma~\ref{lem:hit}. It grows in $P$ while the target's first term
$\sim \sigma' q^{+}/P$ shrinks, and balancing them at
$P \approx \sqrt{\sigma' q^{+}/\delta}$ gives $\sqrt{\sigma' q^{+}\delta}$, which
is precisely the target's second term. That is why $(\mathrm{H1})$ below caps $P$
at the balance point, and it is the only reason.

Note also what the step pays for: on $\{\vx \in I_J\}$ consistency forces
$H^{*}(\vx) = H(\vx)$, the challenge $\Real$ supplies --- so the route pays
$P\delta$ for the event on which $\Dec$'s challenge is \emph{correct}. Consistency
is used nowhere in this section. That is the sense in which the cap is a defect of
the route rather than of the statement, and it is why \S\ref{sec:capping} does not
attack $P\delta$ at all.
\end{remark}

\section{Reducing the hypotheses}\label{sec:reduce}

\begin{lemma}\label{lem:G0}
Let $N \ge 2$. If $8\sqrt{\sigma' q^{+}\delta} \ge 1$ then any conclusion of the
form $\Adv \le 2Bq^{+}/P + 8\sqrt{\sigma' q^{+}\delta} + \gamma$ holds trivially,
its right-hand side being at least $1 \ge \Adv$. Otherwise $q^{+}\delta \le \sigma'$
holds automatically.
\end{lemma}

\begin{proof}
$\Adv$ is a difference of two probabilities, so $\Adv \le 1$; and the right-hand
side is at least $8\sqrt{\sigma' q^{+}\delta}$, every summand being non-negative.
That disposes of the first case. In the second,
$8\sqrt{\sigma' q^{+}\delta} < 1$ gives $\sigma' q^{+}\delta < 1/64$, hence
$q^{+}\delta < 1/(64\sigma')$. Since $\sigma' \ge 2$, $q^{+}\delta < 1/128 < 2
\le \sigma'$.
\end{proof}

\begin{theorem}\label{thm:Epp}
Let $N,M \ge 2$, $\gamma \in (0,1)$, $P \in \mathbb{N}$,
$q \in \mathbb{N} \cup \{0\}$, and let $\mathcal{Y} := \mathcal{Y}^{P,\gamma/2}$ be
the family of Lemma~\ref{lem:P} --- built from $(S_1,S_2,P,\gamma)$ alone, before
$q$, indexed by $(f,\zeta)$, each member a $P$-mixture consistent with its index.
Suppose
\begin{itemize}[leftmargin=1.8em,itemsep=1pt]
\item[$(\mathrm{H1})$] $P \le \sqrt{\sigma' q^{+}/\delta} + 1$, and
\item[$(\mathrm{H2})$] $\min\{\mu'(\min(qM,N^{2})),\ 2\delta\sqrt{M}\}
  \le \sqrt{\sigma' q^{+}\delta}$.
\end{itemize}
Then for every $q$-query challenge-oblivious observer $D$, with no restriction on
challenge resolution,
\[
  \Adv_{\mathcal{Y},D} \le
  \frac{2(\sigma' + \log\gamma^{-1})q^{+}}{P} + 8\sqrt{\sigma' q^{+}\delta}
  + \gamma .
\]
\end{theorem}

\begin{proof}
By Lemma~\ref{lem:G0} we may assume $8\sqrt{\sigma' q^{+}\delta} < 1$, the other
case being trivial. That gives both $\sigma' q^{+}\delta < 1/64$, used in the
sharpened sub-bounds below, and $q^{+}\delta \le \sigma'$. Theorem~\ref{thm:C}
gives, for the same fixed $D$,
\[
  \Adv_{\mathcal{Y},D} \le \varepsilon(D)
  + \frac{2(\sigma' + \log\gamma^{-1})q}{P} + \gamma + P\delta + q\delta,
\]
and the second term is at most the target's first, as $q \le q^{+}$. Four terms
remain.

\emph{The extraction term.} Corollary~\ref{cor:Dpp} applies to $D$ with no
hypothesis and gives
$\varepsilon(D) \le \kappa(q) \le 5\sqrt{\sigma'\delta}
+ \min\{\mu'(\min(qM,N^{2})), 2\delta\sqrt{M}\}$. Here
$5\sqrt{\sigma'\delta} \le 5S$, and the minimum is at most $S$ by
$(\mathrm{H2})$, so $\varepsilon(D) \le 6S$.

\emph{Two sharpened sub-bounds.} Since $\sigma' q^{+}\delta < 1/64$ we have
$q^{+}\delta < 1/(64\sigma')$, so $q^{+}\delta \le \frac{1}{8\sigma'}S$ and, using
$\sigma' \ge 2$,
\[
  q\delta \le q^{+}\delta \le \tfrac{1}{16}S .
\]
Likewise $\delta < 1/(64\sigma' q^{+})$ gives
$\delta \le \frac{1}{8\sigma' q^{+}}S \le \frac{1}{16}S$. Both are worth having:
the crude bounds $q\delta \le S$ and $\delta \le S$ each spend a whole unit of the
target's $8$, and there is no unit to spare.

\emph{The fixed-set term.} With $t_q = \sqrt{\sigma' q^{+}/\delta}$, so that
$t_q\delta = S$, $(\mathrm{H1})$ gives $P \le t_q + 1$ and hence
$P\delta \le t_q\delta + \delta \le S + \frac{1}{16}S = \frac{17}{16}S$.

\emph{The query term.} $q\delta \le \frac{1}{16}S$, above.

Adding, $6 + \frac{17}{16} + \frac{1}{16} = 7\frac18 \le 8$.
\end{proof}

\begin{corollary}[at $q = 0$ there is no restriction on $M$]\label{cor:G2}
Let $N,M \ge 2$, $\gamma \in (0,1)$ and $P \le 3\sqrt{\sigma'/\delta}$. Then for
every $0$-query challenge-oblivious observer and every $M$,
$\Adv_{\mathcal{Y},D} \le 2(\sigma' + \log\gamma^{-1})/P + 8\sqrt{\sigma'\delta}
+ \gamma$.
\end{corollary}

\begin{proof}
$(\mathrm{H2})$ at $q = 0$ reads
$\min\{\mu'(\min(0,N^{2})), 2\delta\sqrt{M}\} = \min\{\mu'(0), 2\delta\sqrt{M}\}$,
and $\mu'(0) = 0$, so $(\mathrm{H2})$ holds for every $M$. In Theorem~\ref{thm:C}
at $q = 0$ the chain is $\Adv \le \varepsilon(D) + \gamma + P\delta$, the two
$q$-carrying terms vanishing. Theorem~\ref{thm:A} gives
$\kappa(0) \le 5\sqrt{\sigma'\delta}$ with no hypothesis on $M$ and none on
resolution --- at $q = 0$, $s = 0$ and every observer has resolution $1$ vacuously
--- so $\varepsilon(D) \le 5\sqrt{\sigma'\delta}$, leaving three units of the
target's $8\sqrt{\sigma'\delta}$ for $P\delta$, and
$P\delta \le 3\sqrt{\sigma'\delta}$ iff $P \le 3\sqrt{\sigma'/\delta}$.
\end{proof}

\begin{remark}[the region for Conjecture~\ref{conj:main} is the intersection over
$q$]\label{rem:A2}
Theorem~\ref{thm:Epp} is a statement \emph{per $q$}: it fixes $P$ and $q$ jointly
and then quantifies over $q$-query observers. Conjecture~\ref{conj:main} is not. It
fixes one family per $(P,\gamma)$ --- before $q$ --- and demands the bound at
\emph{every} $q$. So the conjecture holds at $(P,\gamma)$ exactly when
$(\mathrm{H1})$ and $(\mathrm{H2})$ hold at every $q$, and that is the
intersection, not the region. Both are weakest at small $q$. $(\mathrm{H1})$'s right
side increases in $q$, so it binds at $q = 0$: $P \le \sqrt{\sigma'/\delta} + 1$.
$(\mathrm{H2})$ is vacuous at $q = 0$ and, via its $2\delta\sqrt{M}$ arm, reads
$M \le \sigma' q^{+}/(4\delta)$ for $q \ge 1$, binding at $q = 1$:
$M \le \sigma'/(2\delta)$. This distinction is the single most-mistaken point in
the lineage and is used below.
\end{remark}

\section{The capping construction}\label{sec:capping}

This is the new content, and it does not attack $P\delta$. It observes that one
never has to pay it at the requested $P$.

Given a requested $P \in \mathbb{N}$ and $\gamma \in (0,1)$, put
\[
  \gamma_0 := \max(\gamma,\, N^{-2}),
  \qquad
  P_0 := \min\bigl(P,\ \lceil t_q \rceil\bigr),
  \qquad
  \mathcal{Y} := \mathcal{Y}^{P_0, \gamma_0/2} .
\]

\begin{lemHz}[admissibility]
Every member of $\mathcal{Y}$ is a $P$-mixture consistent with its index;
$\mathcal{Y}$ is indexed by $(f,\zeta)$ and by nothing else; and
Theorem~\ref{thm:Epp} applies to it at $(P_0,\gamma_0)$.
\end{lemHz}

\begin{proof}
Each member is a $P_0$-mixture, so each component has fixed set $I$ with
$|I| \le P_0 \le P$. Definition~\ref{def:bf} requires $|I| \le P$ and nothing more,
so the same witness $(I,a)$ that proves $P_0$-bit-fixing proves $P$-bit-fixing. The
uniformity and independence conditions off $I$ refer to the same unchanged $I$, and
consistency is a property of the fixed values $a$, so both survive verbatim. The
index set is that of $\mathcal{Y}^{P_0,\gamma_0/2}$, namely
$\Fun \times (\bits^{*})^{2}$, as Remark~\ref{rem:index} requires. For
Theorem~\ref{thm:Epp} at $(P_0,\gamma_0)$: $\gamma_0 \in (0,1)$ since $\gamma < 1$
and $N^{-2} \le 1/4$; $(\mathrm{H1})$ reads $P_0 \le t_q + 1$ and holds because
$P_0 \le \lceil t_q \rceil \le t_q + 1$; and $(\mathrm{H2})$ does not mention $P$,
so it is unaffected. Finally $P_0 \ge 1$: $\sigma' \ge 2$, $q^{+} \ge 1$ and
$\delta \le 1$ give $t_q \ge \sqrt2 > 1$, so $\lceil t_q \rceil \ge 2$, and
$P \ge 1$.
\end{proof}

\begin{lemHo}[three slack bounds]
$\log \gamma_0^{-1} \le \log \gamma^{-1}$;
$\ \log \gamma_0^{-1} \le \sigma'$, hence
$\sigma' + \log \gamma_0^{-1} \le 2\sigma'$; and
$N^{-2} \le \tfrac{1}{4}S$.
\end{lemHo}

\begin{proof}
$\gamma_0 \ge \gamma$ gives the first, and it is what keeps $c$ at $2$: without it
the clamp would inflate the leading term rather than the additive one.
$\gamma_0 \ge N^{-2}$ gives $\log \gamma_0^{-1} \le 2\log N \le \sigma'$. For the
third, $\delta \ge 1/N$, $\sigma' \ge 2$ --- which needs $N \ge 2$ --- and
$q^{+} \ge 1$ give $S \ge \sqrt{2/N}$, whence
\[
  N^{-2} \le \tfrac14 S
  \iff N^{-2} \le \tfrac14\sqrt{2/N}
  \iff N^{-3} \le \tfrac12,
\]
true for every $N \ge 2$.
\end{proof}

\begin{thmH}[at fixed $q$, no hypothesis on $P$]
Let $N,M \ge 2$, fix $q \in \mathbb{N} \cup \{0\}$ and suppose $(\mathrm{H2})$. For
every $P \in \mathbb{N}$ and every $\gamma \in (0,1)$, the family $\mathcal{Y}$
above consists of $P$-mixtures consistent with their indices, depends only on
$(S_1,S_2,P,\gamma,q)$, and satisfies, for every $q$-query challenge-oblivious
observer $D$ and with no restriction on challenge resolution,
\[
  \Adv_{\mathcal{Y},D}
  \ \le\ \frac{2(\sigma' + \log\gamma^{-1})\,q^{+}}{P}
  \ +\ 13\,S \ +\ \gamma .
\]
\end{thmH}

\begin{proof}
Lemma~H\textlf{0} licenses Theorem~\ref{thm:Epp} at $(P_0,\gamma_0)$, giving
$\Adv \le 2(\sigma' + \log\gamma_0^{-1})q^{+}/P_0 + 8S + \gamma_0$. Two cases.

\emph{Case (a), $P \le \lceil t_q \rceil$.} Then $P_0 = P$, and by
Lemma~H\textlf{1} $\log\gamma_0^{-1} \le \log\gamma^{-1}$, so the first term is at
most $2Bq^{+}/P$, which is the target's own first term with $c = 2$. For the
additive term, $\gamma_0 \le \gamma + N^{-2} \le \gamma + \tfrac14 S$. Total:
$2Bq^{+}/P + 8\tfrac14 S + \gamma$.

\emph{Case (b), $P > \lceil t_q \rceil$.} Then $P_0 = \lceil t_q \rceil \ge t_q$,
and by Lemma~H\textlf{1} $\sigma' + \log\gamma_0^{-1} \le 2\sigma'$, so
\[
  \frac{2(\sigma' + \log\gamma_0^{-1})q^{+}}{P_0}
  \ \le\ \frac{4\sigma' q^{+}}{t_q}
  \ =\ 4\sigma' q^{+}\sqrt{\frac{\delta}{\sigma' q^{+}}}
  \ =\ 4S .
\]
With $\gamma_0 \le \gamma + \tfrac14 S$ again, the total is at most
$12\tfrac14 S + \gamma \le 13S + \gamma$, and the target's first term is
non-negative.

Both cases are dominated by $2Bq^{+}/P + 13S + \gamma$.
\end{proof}

\begin{remark}[what this is, and is not]\label{rem:relaxation}
This is Conjecture~\ref{conj:main}'s \emph{display} at $c = 2$, $C = 13$, with no
hypothesis on $P$. It is \emph{not} the conjecture, which requires the family to
depend only on $(S_1,S_2,P,\gamma)$; a $q$-dependent family satisfies a relaxation
of it, and Remark~\ref{rem:order} is explicit that the order is part of the
statement. Only one thing changes against the conjecture: the family may depend on
$q$, through $P_0 = \min(P,\lceil t_q \rceil)$ and nothing else. It still may not
depend on $D$, on $\vx$, or on the challenge; it is still one family per
$(P,\gamma,q)$, still indexed by $(f,\zeta)$, still consistent. Against
Theorem~\ref{thm:Epp} the trade is exact: $(\mathrm{H1})$ is deleted outright and
$C$ moves $8 \to 13$, while $c$ does not move. Of the five extra units, four are
case~(b)'s capping and a quarter is the $\gamma$-clamp; three quarters go unused.
\end{remark}

\begin{thmHp}[one $q$-free family, every $P$, at a cost of $\sqrt{q^{+}}$]
Let $N,M \ge 2$ and suppose $(\mathrm{H2})$ holds at \emph{every}
$q \in \mathbb{N} \cup \{0\}$ --- equivalently, by Remark~\ref{rem:A2},
$M \le \sigma'/(2\delta)$, the intersection binding at $q = 1$. Then the family
$\mathcal{Y}' := \mathcal{Y}^{P_0',\gamma_0/2}$ with
$P_0' := \min(P, \lceil t_0 \rceil)$ depends only on $(S_1,S_2,P,\gamma)$ --- not
on $q$ --- consists of $P$-mixtures consistent with their indices, and satisfies,
for every $P \in \mathbb{N}$, every $q \in \mathbb{N} \cup \{0\}$ and every
$q$-query challenge-oblivious observer $D$,
\[
  \Adv_{\mathcal{Y}',D}
  \ \le\ \frac{2(\sigma' + \log\gamma^{-1})\,q^{+}}{P}
  \ +\ \bigl(4\sqrt{q^{+}} + 9\bigr) S \ +\ \gamma .
\]
\end{thmHp}

\begin{proof}
As in Lemma~H\textlf{0}, with $(\mathrm{H1})$ at $P_0'$ holding for \emph{every}
$q$ at once, since $P_0' \le \lceil t_0 \rceil \le t_0 + 1 \le t_q + 1$ by
$t_0 \le t_q$. $(\mathrm{H2})$ is \emph{not} uniform in $q$ and is supplied at every
$q$ by hypothesis. Theorem~\ref{thm:Epp} then applies at $(P_0',\gamma_0)$ for each
$q$. Case (a), $P \le \lceil t_0 \rceil$, is as before. Case (b),
$P > \lceil t_0 \rceil$: now $P_0' \ge t_0 = \sqrt{\sigma'/\delta}$, so
\[
  \frac{4\sigma' q^{+}}{t_0} = 4\sigma' q^{+}\sqrt{\frac{\delta}{\sigma'}}
  = 4q^{+}\sqrt{\sigma'\delta} = 4\sqrt{q^{+}} \cdot S,
\]
using $S = \sqrt{q^{+}}\sqrt{\sigma'\delta}$. Adding $8S$ and
$\gamma_0 \le \gamma + \tfrac14 S$ gives the claim, with a quarter-unit to spare.
\end{proof}

\begin{remark}[the hypothesis is strictly stronger than $(\mathrm{H2})$]
\label{rem:H2every}
An earlier draft wrote ``suppose $(\mathrm{H2})$ holds at $q$'' and then quantified
the conclusion over \emph{every} $q$, leaving $q$ free in the hypothesis and
rebinding it in the conclusion. Three readings, all defective. Read at one $q$, the
hypothesis is \emph{vacuously} satisfiable at $q = 0$, since $\mu'(0) = 0$ makes
$(\mathrm{H2})$ free there, so the theorem would deliver the conjecture for every
$M$ and contradict \S\ref{sec:open}. Read per-$q$, it is true but is then not
Conjecture~\ref{conj:main}, which demands one family good at every $q$. Only the
intersection reading is both true and sufficient, and it was not written down.

The distinction is not idle. A non-vacuous witness: $N = 2^{12}$, $\sigma = 0$ (so
$\sigma' = 24$), $\delta = 2^{-12}$, $M = 2^{16}$. Then
$8\sqrt{\sigma' q^{+}\delta} < 1$ at both $q = 0$ and $q = 1$; $(\mathrm{H2})$
holds at $q = 0$; and at $q = 1$ it fails on both arms, $2\delta\sqrt{M} = 0.125$
and $\mu'(\min(M,N^{2})) = 0.315$, against $S = 0.108$.
\end{remark}

\begin{remark}[relation to the converse construction, and what is not shown]
\label{rem:reduces}
The Contract's own converse remark gives a $q$-free family with
$\Adv \le \kappa(q) + Aq/P_0 + \gamma + P_0\delta + q\delta$ at
$P_0 = \min(P,\lceil\sqrt{A/\delta}\,\rceil)$, $A := \sigma + 2 + \log(1/\gamma)$,
concluding the conjecture ``for $P \le \sqrt{A/\delta}$, and for every $P$ when
$q = O(1)$'', the two directions separated by $\Theta(\sqrt{q^{+}})$ for growing
$q$. Two things must be said carefully.

\emph{It is not the same cap.} $A$ carries a $\log(1/\gamma)$ that
$\sigma' = \sigma + 2\log N$ does not, and has $+2$ where $\sigma'$ has
$2\log N$. The two coincide only at $\gamma = N^{-2}$ and diverge without bound as
$\gamma \downarrow 0$. Theorem~H$'$ is an \emph{analogue} of that construction, run
against Theorem~\ref{thm:Epp} rather than against $\kappa$ directly. It is not a
reproduction of it.

\emph{Only one direction is shown.} It is tempting to conclude that the
$\Theta(\sqrt{q^{+}})$ separation \emph{is} the price of $q$-independence, since
capping at $t_q$ instead of $t_0$ removes it. That claims a necessity nothing here
establishes: no lower bound against $q$-free families appears anywhere, and
\S\ref{sec:open} lists finding such a family --- or showing none exists --- as
open. Moreover the two quantities are different in kind. The converse remark's
$\sqrt{q^{+}}$ is a separation \emph{in the value of $P$} between two directions of
a biconditional; the $\sqrt{q^{+}}$ in Theorem~H$'$ is a \emph{degradation of the
constant $C$ at fixed $P$}. They share an exponent, not a mechanism.
\end{remark}

\section{What $(\mathrm{H2})$ is, and whether applications satisfy it}\label{sec:h2}

$(\mathrm{H2})$ is not a condition about decomposition. It is inherited from routing
the decomposition through extraction, and it enters at exactly one step: the
extraction term of Theorem~\ref{thm:Epp}, where Corollary~\ref{cor:Dpp}'s residue
must fit inside one unit of the budget $S$.

The second arm has a clean reading. By Proposition~\ref{prop:F}'s calculation, a
random function's image distribution $\nu(v) = |f^{-1}(v)|/N^{2}$ is lumpy:
throwing the pair's $\approx 1/\delta^{2}$ effective probability mass into $M$
buckets gives a balls-in-bins deviation $\approx \sqrt{M\delta^{2}} = \delta\sqrt{M}$.
The real challenge is drawn from that lumpy law and the decomposed one is uniform
off the fixed set, so the lumpiness genuinely appears in $\Adv$. Exploiting it,
however, requires knowing \emph{which} values are over-represented, and that means
learning $f$. That is why $q = 0$ is free for every $M$
(Corollary~\ref{cor:G2}), and why $D^{*}$ in Proposition~\ref{prop:F} spends
$N^{2}$ queries.

When $D$ has full challenge resolution --- as a general distinguisher
$D^{H}(y,\vz)$ does, since it may steer its queries on $y$ --- the second arm is
the smaller for $q \ge 1$, $M \ge 4$ by Remark~\ref{rem:whicharm}, and
$(\mathrm{H2})$ becomes, in bits,
\[
  \log M \ \le\ k + \log(\sigma' q^{+}) - 2, \qquad k := \log(1/\delta).
\]
\emph{Do not ask for more output bits than one source's unpredictability}, plus a
logarithmic allowance. The $q$-dependence is only additive-logarithmic: $q$ from
$2^{0}$ to $2^{128}$ moves the allowance by $127$ bits. So the corner where
$(\mathrm{H2})$ binds is not really about query budgets --- it is about asking for a
long output.

\begin{table}[t]
\centering\small
\begin{tabular}{lrrrcc}
\toprule
setting & $k$ & output & allowed & $(\mathrm{H2})$ & bound \\
\midrule
$128$-bit key, $256$-bit-unpredictable sources & $256$ & $128$ & $\le 328$ & yes & $2^{-91}$ \\
$256$-bit key, $256$-bit-unpredictable sources & $256$ & $256$ & $\le 328$ & yes & $2^{-91}$ \\
$128$-bit key, $128$-bit-unpredictable sources & $128$ & $128$ & $\le 200$ & yes & $2^{-27}$ \\
$256$-bit key, $128$-bit-unpredictable sources & $128$ & $256$ & $\le 200$ & \textbf{no} & $2^{-27}$ \\
$512$-bit output, $256$-bit sources & $256$ & $512$ & $\le 328$ & \textbf{no} & $2^{-91}$ \\
$1$\,Mbit stream, $256$-bit sources & $256$ & $2^{20}$ & $\le 328$ & \textbf{no} & $2^{-91}$ \\
\bottomrule
\end{tabular}
\caption{$(\mathrm{H2})$ at application parameters, with $\sigma = 64$,
  $q = 2^{64}$, $\log N = 512$. Computed by
  \texttt{c/0010/campaign/checks/h2-applications.py}.}
\end{table}

The two failing rows differ in kind and only one matters. Row~4 fails
$(\mathrm{H2})$, but the bound there is $2^{-27}$ anyway, so nothing is lost that
one had. Rows~5 and~6 are the real gap: security is still $2^{-91}$ and
$(\mathrm{H2})$ blocks the statement purely for asking for a long output. That is
precisely the sponge and Merkle--Damg\aa{}rd streaming use which motivated wanting a
decomposition route rather than a compression argument in the first place.

\begin{remark}[the evidence for $(\mathrm{H2})$ does not reach the corner that
binds]\label{rem:h2gap}
Proposition~\ref{prop:F}, the only lower bound here, exhibits an observer making
$N^{2}$ queries. But $\kappa$ is monotone in $q$ --- a $q$-query observer may
ignore queries --- so a lower bound at $q = N^{2}$ constrains $\kappa(1)$ not at
all. Meanwhile $(\mathrm{H2})$ is free at $q = N^{2}$, permitting
$M \le \sigma' N^{2}/(4\delta)$, and binds at small $q$. \emph{So nothing here
gives any evidence that $\kappa(q)$ carries $M$-dependence where $(\mathrm{H2})$
actually bites.} The sharp question is therefore narrower than ``can the first arm
of Corollary~\ref{cor:Dpp} be made $M$-free'': does $\kappa(q)$ have \emph{any}
$M$-dependence for small $q$? Exhibit a $q$-query lower bound carrying $M$ for
$q \ll N^{2}$, or prove $\kappa(q) \le O(\sqrt{\sigma'\delta}) + (M\text{-free})$
there.
\end{remark}

\section{What is settled, and what is not}\label{sec:open}

\noindent\textbf{Settled, on the region $(\mathrm{H2})$.} The $P$-axis, at fixed
$q$, in full: no hypothesis on $P$ survives in Theorem~H. And at $q = 0$,
$(\mathrm{H2})$ is vacuous for every $M$ since $\mu'(0) = 0$, so Theorem~H is
unconditional there.

\smallskip
\noindent\textbf{Not settled.}
\begin{itemize}[leftmargin=1.2em,itemsep=2pt]
\item \emph{Conjecture~\ref{conj:main} as stated.} Theorem~H's family depends on
  $q$; Theorem~H$'$'s constant does. Neither is the conjecture, which wants both a
  $q$-free family and an absolute $C$. The honest residue on the $P$-axis is a
  factor $\sqrt{q^{+}}$ in the constant under $q$-free families, on
  $M \le \sigma'/(2\delta)$.
\item \emph{Whether the $\sqrt{q^{+}}$ is necessary.} Only achievability is shown;
  see Remark~\ref{rem:reduces}.
\item \emph{The $M$-corner.} $(\mathrm{H2})$ is carried, not discharged, except at
  $q = 0$. Remark~\ref{rem:h2gap} sharpens the question.
\item \emph{Whether $\delta\sqrt{M}$ is tight off the diagonal.}
  Proposition~\ref{prop:F} licenses nothing off $\delta = 1/N$, and nothing at
  $q < N^{2}$.
\end{itemize}

\smallskip
\noindent\textbf{A redirection.} By Remark~\ref{rem:onestep}, $P\delta$ enters at
one step, and the natural programme is to remove it so that a \emph{large} fixed
set can be afforded. At fixed $q$ that work is unnecessary: one never needs a fixed
set larger than $\lceil t_q \rceil$, so $P\delta$ never needs to exceed
$S + \delta$. Such a programme is aimed solely at the $q$-free statement, where the
cap must sit at $t_0$ and the $\sqrt{q^{+}}$ is real.

\section{Verification status}\label{sec:verification}

\noindent Verification is not uniform across this document, and reading it as
though it were would be the main way to be misled by it.

\smallskip
\noindent\textbf{\S\ref{sec:capping} has had six independent blind referee passes}
in fresh context, spanning two model families: five on the arm itself (lenses: the
definitional step; arithmetic and constants; quantifier order; case coverage and
degenerate instances; citation fidelity) and one on three proposed simplifications.
All five returned \textsc{defects}. Lemma~H\textlf{0} was attacked by four referees
independently, each trying to construct a reading of Definition~\ref{def:bf} on
which it fails; none could. The arithmetic referee re-derived every inequality by
hand, reproduced the machine check's grid sizes exactly including collision counts,
found \emph{no computational error}, and confirmed that $c = 2$ and $C = 13$ follow
from the displayed inequalities rather than from the grid. Seven class-(A) and
class-(C) findings were upheld and repaired, every one a claim \emph{about} the
results rather than a result: the malformed hypothesis of Theorem~H$'$
(Remark~\ref{rem:H2every}); an identification of the $\sqrt{q^{+}}$ with the
converse remark's separation (Remark~\ref{rem:reduces}); the unstated standing
hypothesis $N,M \ge 2$; a machine check described as reaching $q \le 10^{7}$ whose
vacuity filter admitted nothing above $q^{+} = 410$; and summary lines scoped to a
stronger statement than the body proves. The revision carrying those repairs has
itself had \emph{no} pass.

\smallskip
\noindent\textbf{\S\ref{sec:reduce} has had two passes}, both \textsc{defects},
both repaired; the current form has had none.

\smallskip
\noindent\textbf{\S\S\ref{sec:extraction}--\ref{sec:lower} have had none.} This is
the weakest link and it is load-bearing. Corollary~\ref{cor:Dpp} supplies the
extraction term of Theorem~\ref{thm:Epp}, on which everything in
\S\ref{sec:capping} rests, and it sits in an arm with zero verification passes.
Worse, the one referee objection ever formed about its constant was \emph{overruled}
rather than upheld, and never re-tested. A reader who wants one thing checked
should check Corollary~\ref{cor:Dpp}.

\smallskip
\noindent\textbf{\S\S\ref{sec:structural}--\ref{sec:revealing} and
\S\ref{sec:presampling} have had one pass}, verdict \textsc{clean}, with six
class-(C) findings that were never adjudicated.

\smallskip
\noindent\textbf{The external inputs have had none.}
Assumptions~\ref{ass:cdgs2}--\ref{ass:binom} are quoted from source cards
transcribed from local copies, and no pass has ever checked a card against the
paper it summarises. Three consequences of Assumption~\ref{ass:cdgs2} used in
\S\ref{sec:presampling} are not stated in the published claim in that form.

\smallskip
\noindent\textbf{Lemmas~\ref{lem:hit} and~\ref{lem:query} are unproved here and
unrefereed anywhere.} They are the Contract's, they are load-bearing inside
Theorem~\ref{thm:C} at exactly the step where the $P$-cap enters, and no
verification pass has examined them, the Contract having been handed to every
referee as the yardstick rather than as an object of review.

\smallskip
\noindent\textbf{No prior-art search has ever been run on this campaign}, so the
novelty of Theorem~\ref{thm:D}, Corollary~\ref{cor:Dp} and
Proposition~\ref{prop:F} --- all unconditional --- is unchecked. And no
counterexample search has been run against any current artifact.

\end{document}
